\begin{thebibliography}{99}
\bibitem{bib-bib1} 
\par\noindent    Peter Aczel.  The Russell–Prawitz modality.  Mathematical. Structures in Comp. Sci., 11(4):541–554, August 2001.  doi:10.1017/S0960129501003309.  

\bibitem{bib-bib2} 
\par\noindent    Michael Barr.  Toposes without points.  Journal of Pure and Applied Algebra, 5(3):265–280, 1974.  doi:10.1016/0022-4049(74)90037-1.  

\bibitem{bib-bib3} 
\par\noindent    Chad E. Brown and Christine Rizkallah.  Glivenko and Kuroda for simple type theory.  The Journal of Symbolic Logic, 79(2):485–495, 2014.  doi:10.1017/JSL.2013.10.  

\bibitem{bib-bib4} 
\par\noindent    Alonzo Church.  A Formulation of the Simple Theory of Types.  The Journal of Symbolic Logic, 5(2):56–68, 1940.  doi:10.2307/2266170.  

\bibitem{bib-bib5} 
\par\noindent    Almudena Colacito, Dick de Jongh, and Ana Lucia Vargas.  Subminimal negation.  Soft Computing, 21:165–174, 2016.  doi:10.1007/s00500-016-2391-8.  

\bibitem{bib-bib6} 
\par\noindent    Martín Escardó and Paulo Oliva.  The Peirce translation.  Annals of Pure and Applied Logic, 163(6):681–692, 2012.  Computability in Europe 2010.  doi:10.1016/j.apal.2011.11.002.  

\bibitem{bib-bib7} 
\par\noindent    Matt Fairtlough and Michael Mendler.  Propositional Lax Logic.  Information and Computation, 137:1–33, 1997.  doi:10.1006/inco.1997.2627.  

\bibitem{bib-bib8} 
\par\noindent    Gilda Ferreira and Paulo Oliva.  On the Relation Between Various Negative Translations.  In Ulrich Berger, Hannes Diener, Peter Schuster, and Monika Seisenberger, editors, Logic, Construction, Computation, pages 227–258. De Gruyter, Berlin, Boston, 2012.  doi:10.1515/9783110324921.227.  

\bibitem{bib-bib9} 
\par\noindent    Harvey Friedman.  Classically and intuitionistically provably recursive functions.  In Gert H. Müller and Dana S. Scott, editors, Higher Set Theory, pages 21–27, Berlin, Heidelberg, 1978. Springer Berlin Heidelberg.  

\bibitem{bib-bib10} 
\par\noindent    Gerhard Gentzen.  Die Widerspruchsfreiheit der Reinen Zahlentheorie.  Mathematische Annalen, 112:493–565, 1936.  

\bibitem{bib-bib11} 
\par\noindent    Valery Glivenko.  Sur quelques points de la logique de M. Brouwer.  Bulletins de la classe des sciences, 15:183–188, 1928.  

\bibitem{bib-bib12} 
\par\noindent    Kurt Gödel.  Zur intuitionistischen Arithmetik und Zahlentheorie.  Ergebnisse eines Mathematischen Kolloquiums, 4:34–38, 1933.  

\bibitem{bib-bib13} 
\par\noindent    Andrey Nikolaevich Kolmogorov.  O principe tertium non datur.  Matematicheskiĭ Sbornik, 32:646–667, 1925.  

\bibitem{bib-bib14} 
\par\noindent    Sigekatu Kuroda.  Intuitionistische Untersuchungen der formalistischen Logik.  Nagoya Mathematical Journal, 2:35–47, 1951.  doi:10.1017/S0027763000010023.  

\bibitem{bib-bib15} 
\par\noindent    Eugenio Moggi.  Notions of computation and monads.  Information and Computation, 93(1):55–92, 1991.  Selections from 1989 IEEE Symposium on Logic in Computer Science.  doi:10.1016/0890-5401(91)90052-4.  

\bibitem{bib-bib16} 
\par\noindent    Sara Negri.  Contraction-free sequent calculi for geometric theories with an application to Barr’s theorem.  Archive for Mathematical Logic, 42:389–401, May 2003.  doi:10.1007/s001530100124.  

\bibitem{bib-bib17} 
\par\noindent    Sara Negri.  Geometric rules in infinitary logic.  In Ofer Arieli and Anna Zamansky, editors, Arnon Avron on Semantics and Proof Theory of Non-Classical Logics, pages 265–293. Springer International Publishing, Cham, 2021.  doi:10.1007/978-3-030-71258-7\_12.  

\bibitem{bib-bib18} 
\par\noindent    Erik Palmgren.  An Intuitionistic Axiomatisation of Real Closed Fields.  Mathematical Logic Quarterly, 48(2):297–299, 2002.  doi:10.1002/1521-3870(200202)48:2\textbackslash{}\%3C297::AID-MALQ297\textbackslash{}\%3E3.0.CO;2-G.  

\bibitem{bib-bib19} 
\par\noindent    Luiz Carlos Pereira and Edward Hermann Haeusler.  On constructive fragments of classical logic.  In Heinrich Wansing, editor, Dag Prawitz on Proofs and Meaning, pages 281–292. Springer International Publishing, Cham, 2015.  doi:10.1007/978-3-319-11041-7\_12.  

\bibitem{bib-bib20} 
\par\noindent    Michael Rathjen.  Remarks on Barr’s Theorem: Proofs in Geometric Theories.  In Dieter Probst and Peter Schuster, editors, Concepts of Proof in Mathematics, Philosophy, and Computer Science, pages 347–374. De Gruyter, Berlin, Boston, 2016.  doi:10.1515/9781501502620-019.  

\bibitem{bib-bib21} 
\par\noindent    Thomas Traversié.  Kuroda’s Translation for Higher-Order Logic.  Submitted, April 2024.  doi:10.48550/arXiv.2404.19503.  

\bibitem{bib-bib22} 
\par\noindent    Benno van den Berg.  A Kuroda-style j-translation.  Archive for Mathematical Logic, 58(5-6):627–634, August 2019.  doi:10.1007/s00153-018-0656-x.  

\end{thebibliography}
